\section{Introduction}
\label{sec:intro}

\IEEEPARstart{S}{elf}-supervised representation learning has transformed
language and vision, and is now reshaping electroencephalography (EEG)
analysis. Large EEG foundation models---LaBraM~\cite{labram2024},
EEGPT~\cite{eegpt2024}, NeuroLM~\cite{neurolm2025},
BIOT~\cite{biot2023}---pretrain on thousands of hours of EEG with 5--10\,M$+$
parameters on multi-GPU clusters, and transfer via fine-tuning or linear
probing to clinical and brain--computer interface (BCI) tasks. These
methods establish that scale helps, but their compute and parameter budgets
are incompatible with on-device deployment in wearable EEG hardware (Muse,
OpenBCI, Neurosity Crown, Smartbuds-class earbud devices) where latency,
energy, and memory budgets are measured in milliseconds, milliwatts, and
megabytes respectively.

\emph{When does compact SSL pretraining actually help on EEG, and how should
practitioners spend a limited compute budget?} The contemporary literature
provides partial answers. LeJEPA~\cite{lejepa2025} introduced the Sketched
Isotropic Gaussian Regularization (SIGReg) anti-collapse objective and a
small JEPA-style world model demonstrating $\sim$15\,M parameters trained on
a single GPU could be competitive at world-modelling. Concurrent work
Laya~\cite{laya2026} applied LeJEPA to EEG at foundation scale (29{,}000
hours of EEG, 20{,}940 subjects, two L40S GPUs), establishing that latent
prediction outperforms reconstruction-based pretraining on clinical EEG
benchmarks under \emph{cross-dataset evaluation}: EEGMMIDB and BCI Competition~IV
are held out from Laya's pretraining and pooled into multi-dataset downstream
benchmarks, on which no SSL foundation model exceeds 0.51 balanced accuracy
on left-versus-right motor imagery~\cite[Tab.~1]{laya2026}. The complementary
\emph{in-distribution single-corpus regime}---relevant to many applied EEG
settings where a research lab pretrains and evaluates on a focused corpus
realistic to its experimental design---is not addressed. Neither paper
characterizes the small-scale scaling regime, the SIGReg-weight optimum, or
the cross-dataset channel-transfer behaviour that practitioners face when
deploying a single pretrained encoder across heterogeneous downstream
recording rigs.

This paper fills that gap. We present \textbf{EEG-LeJEPA}, a 2.86\,M-parameter
encoder-plus-causal-predictor architecture trained with
SIGReg~\cite{lejepa2025} on 50 subjects of EEGMMIDB motor-imagery data---a
deliberately compact-scale recipe runnable on a single consumer GPU in 10
minutes. Across six (data, compute) operating points we map the scaling law,
we explicitly characterize the SIGReg-weight optimum, and we compare two
strategies for cross-dataset transfer (channel padding vs.\ channel-matched
pretraining). We compare against (i) a randomly-initialized identical
architecture as the no-pretraining control and (ii) the same architecture
trained end-to-end with supervised cross-entropy as the no-SSL control.

\subsection*{Contributions}

\begin{enumerate}
\item A compact (2.86\,M parameter) SIGReg-pretrained EEG encoder that
achieves 68.6\,\% accuracy / 0.751 AUC on EEGMMIDB rest-versus-activity
and 62.1\,\% / 0.690 AUC on EEGMMIDB left-vs-right motor imagery via a
frozen-encoder linear probe on encoder mean-pooled features
(leave-one-subject-out with subjects~1--20 strictly held out of
pretraining; both gains over random significant, $q<0.002$), with the
full pretraining$+$evaluation pipeline running on a single consumer GPU in
under 15\,minutes at total compute cost under \$1.

\item A scaling-law characterization across six (data, compute) operating
points (3--89 subjects, 200--30k training steps) under a strictly
leakage-free split, showing the compact model's best-source AUC rises and
plateaus near 0.76. A 2.4$\times$ larger model (7\,M parameters) trained on
the same data is modestly but consistently better ($+1.3$ to $+2.0$\,pp on
encoder mean-pooled features, seed-stable across three seeds), indicating the
plateau is a capacity limit; the compact model is retained as a favorable
accuracy/size trade-off for edge deployment. An isolation run shows the
earlier ``larger model regresses'' result was an artifact of too few SIGReg
projections at the wider embedding, not of capacity.

\item A SIGReg-weight ablation across $\lambda \in \{0, 0.1, 0.3, 1.0,
3.0, 10.0\}$ showing a sharp U-shape with optimum at $\lambda=1.0$ and
\emph{complete representational collapse} at $\lambda=0$ (probe accuracy
regresses to exactly chance), providing an empirical validation of the
SIGReg anti-collapse mechanism that complements the theoretical analysis
of~\cite{lejepa2025}.

\item A counterintuitive cross-dataset transfer finding: pretraining on
the channel intersection with the target dataset
\emph{underperforms} channel-padded transfer from the full-channel
source. We attribute this to the channel-rich pretrained encoder learning
spatially-redundant representations that gracefully degrade under
input-channel reduction. For practitioners deploying a single pretrained
EEG encoder across heterogeneous recording rigs, this argues for
pretraining at the richest available montage and padding at inference.

\item A direct comparison to supervised-from-scratch training of the same
architecture: our SSL$+$linear-probe pipeline ties end-to-end supervised
training on left-vs-right MI ($-0.2$\,pp, n.s.) and significantly exceeds
it on rest-vs-activity ($+8.4$\,pp, $p<0.001$), while a single frozen
encoder serves both downstream tasks at sub-second per-task probing cost.

\item The complete reproducibility pipeline: 39 unit tests across six
test modules, the trained checkpoints used in this paper (six
$\lambda$-ablation checkpoints plus the scaling, capacity, and
channel-transfer checkpoints), per-fold CSV outputs for every reported
table, and end-to-end training$+$probe runtime under 15 minutes on a
single RTX-class GPU. Code, configurations and checkpoints released
at \url{https://github.com/lifeinpassion/eeg-lejepa}.
\end{enumerate}

The remainder of the paper is organized as follows.
\cref{sec:related} surveys EEG foundation models and SSL anti-collapse
mechanisms with explicit relation to concurrent work
\cite{laya2026}. \cref{sec:method} presents the EEG-LeJEPA architecture
and SIGReg objective. \cref{sec:experiments} details the empirical setup
and downstream evaluation protocols. \cref{sec:results} presents the
scaling, ablation and transfer results. \cref{sec:discussion} situates
our findings against published EEG foundation models and outlines
limitations and future work. \cref{sec:conclusion} concludes.
